\section{Solución Propuesta}\label{sec:sol_prop}

La solución propuesta consiste en extender experimentalmente la implementación existente de \texttt{ApplicativeDo} en el compilador \textsf{GHC}, con el objetivo de que el pipeline actual pueda considerar reordenamientos semánticamente válidos de \textit{statements} dentro de la \textit{do-notation} cuando se trabaja con mónadas conmutativas. En términos generales, la memoria no propone construir un compilador nuevo ni un lenguaje independiente, sino intervenir una transformación ya presente en \textsf{GHC}, tomando como base el algoritmo de \textit{Rearrangement} y \textit{Desugaring} propuesto por Marlow et al. \cite{MarlowPJKM2016}. La idea central es ampliar este pipeline para que no se limite al orden sintáctico original de los 	\textit{statements}, sino que pueda evaluar permutaciones válidas y seleccionar el plan de ejecución de menor costo según el modelo adoptado, sin alterar la semántica del programa.

\newpage

Para llevar a cabo esta propuesta, se trabajará sobre un repositorio experimental que ya incorpora un entorno reproducible mediante el uso de un \textit{devcontainer} y un submódulo con el código fuente del compilador \textsf{GHC}. Sobre este entorno, será necesario disponer de la cadena de herramientas base del compilador, incluyendo una instalación funcional de \textsf{GHC} para bootstrap, \textit{cabal-install} para la gestión y construcción de componentes \textsf{Haskell}, y las herramientas \textit{happy} y \textit{alex}, utilizadas por \textsf{GHC} para generar analizadores sintácticos y léxicos. Adicionalmente, el proceso de compilación del compilador requiere herramientas estándar del sistema, tales como un compilador de \texttt{C} y utilidades de configuración, las cuales se consideran encapsuladas dentro del entorno de desarrollo reproducible preparado para esta memoria.

La primera etapa de la solución consistirá en diseñar e implementar un algoritmo capaz de generar todas las permutaciones válidas de \textit{statements} dentro de una mónada conmutativa, respetando las restricciones \textit{def-use} y preservando la semántica del programa. Esta etapa corresponde al componente nuevo de la propuesta y busca ampliar el espacio de configuraciones candidatas que luego serán evaluadas por el algoritmo base de Marlow et al. En particular, el propósito no es reemplazar el pipeline original, sino producir un conjunto de órdenes sintácticos alternativos que sean semánticamente admisibles en el contexto probabilístico y, por tanto, susceptibles de ser ingresados posteriormente al proceso de \textit{Rearrangement} y \textit{Desugaring}.

Una vez generadas las permutaciones válidas, estas se integrarán al algoritmo propuesto por Marlow et al., de modo que cada una produzca un plan de ejecución candidato. Luego, dichos planes se compararán usando el modelo de costos del artículo base, seleccionando el de menor costo. De esta manera, el algoritmo original no se reemplaza, sino que pasa a operar sobre un conjunto más amplio de configuraciones sintácticamente posibles y semánticamente válidas. Sobre esta base, la extensión experimental en \textsf{GHC} consistirá en incorporar el reordenamiento de \textit{statements} dentro del pipeline actual de \texttt{ApplicativeDo}. Así, la propuesta se sitúa en una modificación controlada del compilador, y no en una transformación externa al mismo. En consecuencia, el objetivo es que el propio compilador pueda considerar órdenes alternativos de evaluación cuando estos sean válidos, y escoger aquel que produzca el plan de ejecución más conveniente según el criterio de costo adoptado.

\newpage



